\documentclass{article}
\usepackage[utf8]{inputenc}
\usepackage[italian]{babel}
\usepackage{fancyhdr}
\usepackage{geometry}
\usepackage{graphicx}
\usepackage{wrapfig}
\usepackage{gensymb}
\usepackage{appendix}
 \geometry{
 a4paper,
 total={180mm,257mm},
 left=15mm,
 top=20mm,
 }
 \usepackage{multicol}
 \usepackage{color}
 \setlength{\columnseprule}{0.4pt} %Separator ruler width
\def\columnseprulecolor{\color{black}} %Separator ruler colour
\usepackage[
backend=biber,
style=alphabetic,
sorting=ynt
]{biblatex}
\addbibresource{name.bib}
\pagestyle{fancy}
\fancyhf{}
\rhead{Corso di Laurea in Fisica - A.A. 2019/20 \\ Università degli Studi di Torino}
\lhead{Esperimentazioni II (corso B) \\ Modulo I - Prof. Roberto Bellan}
\rfoot{\thepage}
\begin{document}
% inserire titolo dell'esperienza

\noindent
\textbf{\LARGE Strumenti di misura, caratteristica I-V di una lampadina ed interruttori}
\\
\\
\\
\textbf{\large Thomas Robiglio\textsuperscript{1}, Lorenzo Stefanuto e Sergio G. Tamburi}\\
\textsuperscript{1} Thomas Robiglio, e-mail: \verb|thomas.robiglio@edu.unito.it|
\begin{table}[h!]
\begin{tabular}{|l|l|}
\hline
Inviato  & 03/03/2020              \\ \hline \hline
Ricevuto & ..................... \\ \hline \hline
Valutato & ..................... \\ \hline
\end{tabular}
\end{table}

% inserire obbiettivi di misura
\hrule
\begin{center}
\textbf{\underline{\large OBIETTIVI DELL'ESPERIENZA}}
\end{center}
L'esperienza svolta è articolata in 3 parti distinte, ciascuna con obiettivi differenti.\\
La prima parte dell'esperienza mira a determinare la resistenza interna di un tester digitale (\textsc{Amprobe 37xr-a}) e di un tester analogico (\textsc{Supertester 680r}), impiegati sia come voltmetro che come amperometro, andando a calcolare anche la potenza consumata dagli strumenti per fornire le misure.\\
La seconda parte dell'esperienza è volta a determinare, effettuando misure di corrente e tensione, la caratteristica \(I-V\) di una lampadina e a verificare la validità della legge di corpo nero per un filamento di tungsteno riscaldato.\\
La terza parte dell'esperienza - puramente operativa - consiste nell'implementare, utilizzando due switch, un generatore di tensione e una lampadina, un circuito \textsc{And}, un circuito \textsc{Or} e un circuito \textit{two-way switching}.
\\
\hrule
\begin{multicols}{2}
\noindent
\underline{NOTA}: Per le procedure sperimentali e la simbologia si è fatto riferimento alle dispense fornite dal docente.\\
\\
Per la prima parte dell'esperienza, in riferimento a quanto presente nelle dispense, si è impiegata una resistenza da \((487\pm3)k\Omega\) misurata con \textsc{Agilent u1731a}.  
\section{Elaborazione dei dati}
\subsection{Multimetri impiegati come voltmetri}
Per determinare la resistenza interna dei multimetri impiegati come voltmetri, dalla formula per il partitore di tensione \cite{boyestad} si ha: 
\begin{equation}
    \frac{V_{AG}}{R+R_i}=V_{BG} \Rightarrow R_i=\frac{V_{BG}\cdot R}{V_{AB}-V_{BG}}
\end{equation}
La potenza consumata si determina a partire dalla legge di Ohm \cite{boylestad}:
\begin{equation}
    P=\frac{V_{BG}^2}{R_i}=\frac{V_{AB}-V_{BG}}{R}\cdot V_{BG}
\end{equation}
\subsection{Multimetri impiegati come amperometri}
La resistenza interna e la potenza dissipata si ricavano direttamente dalla legge di Ohm.
\subsection{Potenza dissipata per una lampadina a filamento di tungsteno}
Dalle letteratura (\cite{agrawal} e \cite{clauss}) sono note, per una lampadina a filamento di tungsteno, le relazioni:
\begin{equation}
    P=k(T-T_0)+e\sigma A_s(T^4 - T_0^4)
\end{equation}
\begin{equation}
    R=R_0 (\frac{T}{T_0})^b \Rightarrow T=\beta R^{\gamma}
\end{equation}
Note (3) e (4) è possibile operare alcune apprrossimazioni per descrivere i dati raccolti.\\  
In prima approssimazione (\textbf{ipotesi 1}) si trascura sia il termine conduttivo e convettivo della potenza dissipata sia il termine \(e\sigma A_s T_0^4\) - si ipotizza infatti che la temperatura iniziale della lampadine \(T_0\) sia trascurabile rispetto alle temperature raggiunte dal filamento (\(>  2000 \celsius\)); si ha quindi una relazione del tipo:
\begin{equation}
    P=mR^q
\end{equation}
L'altra ipotesi (\textbf{ipotesi 2}) tiene conto di tutti i termini presenti in (3):
\begin{equation}
    P=mR^q +nR^{\frac{q}{4}}+K
\end{equation}
Si noti che in (6) il termine \(K\) deve assumere valore negativo.\\
Si procede quindi effettuando delle regressioni sui dati sperimentali (avendo ricavato dai dati raccolti i valori di \(R\) e \(P\) con la legge di Ohm) con le forme funzionali delle 2 ipotesi. 
\subsection{Caratteristica I-V di una lampadina a filamento di tungsteno}
Dalla relazione (7) si può ricavare (ricordando che \(P=VI\) e \(R=\frac{V}{I}\)) la relazione tra la corrente e la tensione ai capi della lampadina:
\begin{equation}
VI=m(\frac{V}{I})^q \Rightarrow I=sV^{p}   
\end{equation}
dove \(s=\sqrt[q+1]{m}\) e \(p=\frac{q-1}{q+1}\).\\
Si esegue quindi una regressione sui dati di \(I\) e \(V\) raccolti per verificare sperimentalmente la coerenza tra l'ipotesi di corpo nero e la caratteristica \(I-V\) della lampadina.


\section{Risultati}
In riferimento a \textsection 1.1 i valori ottenuti per lo strumento \textsc{Supertester 680r} con fondo scala a \(10V\) sono:
\begin{itemize}
    \item \(R_i=(180 \pm 13) k\Omega\)
    \item \(P=(2,2 \pm 0,3)\cdot 10^{-5}W\)
\end{itemize}
Per lo strumento \textsc{Amprobe 37xr-a}:
\begin{itemize}
    \item \(R_i=(9,2 \pm 0,4) M\Omega\)
    \item \(P=(6,2 \pm 0,3)\cdot 10^{-6}W\)
\end{itemize}
In riferimento a \textsection 1.2 i valori ottenuti per lo strumento \textsc{Supertester 680r} con fonda scala a \(5mA\) sono:
\begin{itemize}
    \item \(R_i=(63 \pm 2) \Omega\)
    \item \(P=(18,3 \pm 0,6)\cdot 10^{-5}W\)
\end{itemize}
Per lo strumento \textsc{Amprobe 37xr-a}:
\begin{itemize}
    \item \(R_i=(11,1 \pm 0,2) \Omega\)
    \item \(P=(32,8 \pm 0,7)\cdot 10^{-6}W\)
\end{itemize}
In riferimento a \textsection 1.3 solo l'\textbf{ipotesi 2} risulta significativa  e coerente con le previsioni teoriche sull'intero set di misure, con i parametri:
\begin{center}
\begin{tabular}{|l|l|}
\hline
parametro  \(m\) & \((6,3\pm1,4) \cdot 10^{-6}\frac{W}{\Omega^q}\) \\ \hline
parametro \(n\)  & \((0,55\pm0,02) \frac{W}{\Omega^{\frac{q}{4}}}\)          \\ \hline
parametro \(q\)  & \((3,96\pm0,05\))                             \\ \hline
parametro \(K\)  & \((-5,2\pm0,5)W\)                           \\ \hline
\end{tabular}
\end{center}
Con cui si ha il valore:
\[\gamma=\frac{q}{4}=(0,989\pm 0,012)\]
\noindent
L'\textbf{ipotesi 1} risulta significativa considerando le ultime 19 coppie di valori; con i parametri:
\begin{center}
    \begin{tabular}{|l|l|}
\hline
parametro \(m\) & \((2,8\pm0,4) \cdot 10^{-5} \frac{W}{\Omega^q}\) \\ \hline
parametro \(q\) & \((3,65\pm0,03\))                                \\ \hline
\end{tabular}
\end{center}
Con cui si ha il valore:
\[\gamma=(0,913\pm 0,008)\]
In riferimento a \textsection 1.4 i parametri ottenuti dalla regressione (eseguita sulle ultime 19 coppie di valori, coerentemente con quanto visto sopra) sono: 
\begin{center}
    \begin{tabular}{|l|l|}
\hline
parametro \(s\) & \((23,76\pm0,17) \frac{A}{V^p}\) \\ \hline
parametro \(p\) & \((0,571\pm0,004\))             \\ \hline
\end{tabular}
\end{center}
É quindi possibile ricavare il valore dell'esponente \(q\):
\[p=\frac{q-1}{q+1} \Rightarrow q=\frac{1+p}{1-p}=(3,66\pm0,02)\]
Che da il valore:
\[\gamma=(0,916\pm 0,005)\]
\section{Osservazioni e conclusioni}
I risultati della prima parte dell'esperienza appaiono apprezzabili e coerenti con le previsioni teoriche \cite{boylestad} per i valori delle resistenze interne di voltmetri e amperometri.\\
L'analisi della seconda parte dell'esperienza ha mostrato che la legge di corpo nero descrive in maniera soddisfacente il comportamento \(P(R)\) di una lampadina solo a temperature sufficientemente elevate; considerando l'intero set di dati i termini di potenza consumata per via convettiva e conduttiva non risultano trascurabili.\\
Sul set di dati ristretto l'ipotesi di corpo nero si è dimostrata consistente con la caratteristica \(I-V\) della lampadina (compatibilità tra i valori del parametro \(q\) ottenuti).\\
I valori sperimentali di \(\gamma\) determinati nelle due differenti ipotesi non sono risultati compatibili tra loro; si ipotizza che tale risultato sia legato a un comportamento diverso da quello descritto da (4) per basse temperature. Inoltre tali valori non si sono mostrati significativamente consistenti con i valori teorici (nonchè privi di errore) presenti in letteratura (\cite{gluck}, \cite{clauss} e \cite{agrawal}), si evidenzia quindi la necessità di approfondire le ricerche per evidenziare eventuali fonti di errore non considerate.\\ 
La terza parte dell'esperienza è stata svolta con successo; sono state verificate le tavole di verità e i comportamenti dei tre circuiti.
\end{multicols}
\medskip
\noindent
\textit{Si presentano in appendice i grafuci, i parametri dei test di verifica delle bontà dei fit per i risultati ottenuti e le modellizzazioni dei circuiti implementati nella terza parte dell'esperienza.}

\medskip
 
\printbibliography

\newpage
\begin{center}
\textbf{\underline{\Large APPENDICE}}
\end{center}
In riferimento a \textsection 1.3, per l'\textbf{ipotesi 2}:
\begin{figure}[h!]
\caption{Grafico \(P(R)\) per l'\textbf{ipotesi 2}}
\centering
\includegraphics[width=0.7\textwidth]{noapprossimazioni.png}
\end{figure}
\begin{tabular}{c|c}
\end{tabular}
\begin{center}
\begin{tabular}{|c|c|}
\hline
\(\chi^2\)      & 6,55   \\ \hline \hline
Gradi di libertà & 20     \\ \hline \hline
Probabilità      & 0,9979 \\ \hline
\end{tabular}
\end{center}
\medskip
In riferimento a \textsection 1.3, per l'\textbf{ipotesi 1}, dalla regressione eseguita sulle ultime 19 coppie di valori:
\begin{figure}[h!]
\caption{Grafico \(P(R)\) per l'\textbf{ipotesi 1} - Radiazione di corpo nero}
\centering
\includegraphics[width=0.7\textwidth]{corponero.png}
\end{figure}
\begin{tabular}{c|c}
\end{tabular}
\begin{center}
\begin{tabular}{|c|c|}
\hline
\(\chi^2\)      & 11,15   \\ \hline \hline
Gradi di libertà & 16     \\ \hline \hline
Probabilità      & 0,8002 \\ \hline
\end{tabular}
\end{center}
\newpage
In riferimento a \textsection 1.4 dalla regressione eseguita sulle ultime 19 coppie di valori:
\begin{figure}[h!]
\caption{Grafico \(I(V)\) - Caratteristica I-V di una lampadina}
\centering
\includegraphics[width=0.7\textwidth]{corrente.png}
\end{figure}
\begin{tabular}{c|c}
\end{tabular}
\begin{center}
\begin{tabular}{|c|c|}
\hline
\(\chi^2\)      & 8,21   \\ \hline \hline
Gradi di libertà & 16     \\ \hline \hline
Probabilità      & 0,9423 \\ \hline
\end{tabular}
\end{center}
\medskip
Per la terza parte dell'esperienza i 3 circuiti sono stati implementati come mostrato nelle figure.
\begin{figure}[h!]
\caption{Circuito \textsc{And}}
\centering
\includegraphics[width=0.3\textwidth]{AND.png}
\end{figure}
\begin{figure}[h!]
\caption{Circuito \textsc{Or}}
\centering
\includegraphics[width=0.3\textwidth]{OR.png}
\end{figure}
\begin{figure}[h!]
\caption{Circuito \textit{two-way switching}}
\centering
\includegraphics[width=0.3\textwidth]{switch.png}
\end{figure}

\end{document}
